% =====================================================================
\section{The recurrent layer}\label{sec:chain}
% =====================================================================

The cascade-pipeline substrate (\S\ref{sec:substrate}) reproduces
cortical-avalanche statistics matching real EEG (\S\ref{sec:results}).
To test high-level signatures — NREM-N3 variance collapse, propofol
regime-switching, propofol $\Phi$ collapse — we add a recurrent
self-model layer above the substrate. The layer adds one
graph-pinned nonlinearity, one condition-dependent self-injection
coupling $\eta$, and four trajectory observables. No shape parameter
is fit to any neural dataset.

This section is method, not metaphysics. We do not claim the
recurrent layer ``is'' consciousness; we report which numerical
observables on the layer's trajectory match published biological
signatures in~\S\ref{sec:results}.

\subsection{The recurrent loop}

Implementation: \texttt{kernel/self\_model\_stream.py:SelfModelLoop}.
At each tick $t$ the substrate state evolves as
\begin{align}
f_{\mathrm{total}}(t) &= f_{\mathrm{ext}}(t) + \eta\cdot f_{\mathrm{self}}(\mathrm{snap}_{t-1}, \psi_{t-1}), \\
\psi_{t} &= \Cph^{-1}\, f_{\mathrm{total}}(t), \\
\psi^{\mathrm{thr}}_{t} &= \mathrm{bounded\_topk}(\psi_{t}, k=12), \\
\mathrm{state}_{t} &= \mathrm{decay}\cdot\mathrm{state}_{t-1} + (1-\mathrm{decay})\cdot \psi^{\mathrm{thr}}_{t}, \\
\mathrm{snap}_{t} &= \mathrm{bind\_phenomenal\_field}(\psi_{t}, \texttt{profile=K\_7\_only}),
\end{align}
with $\mathrm{decay}=0.95$ (state EMA factor) and $\eta$ the only
condition-dependent architectural parameter. $f_{\mathrm{self}}$ maps
the prior phenomenal snapshot to a directional source weighted by
ignition $\times$ intensity (cosine direction alignment with the
prior snapshot). The substrate response operator $\Cph$ is unchanged
across all conditions.

Conditions:
\begin{itemize}\itemsep=2pt
\item $\eta = 0.20$ for WAKE and RECOVERY (active recurrent self-loop);
\item $\eta = 0.05$ for SLEEP\_N3 (attenuated self-loop);
\item $\eta = 0.00$ for PROPOFOL (broken recurrence; residual cortex
  preserved as background drive).
\end{itemize}

\subsection{The graph-pinned nonlinearity}

\textbf{$\mathrm{bounded\_topk}(\psi, k=12)$.} This is the load-bearing
nonlinearity, implemented in
\texttt{kernel/lyapunov\_selector.py:bounded\_topk}: zero all but the
top-$12$ vertex amplitudes (by absolute value), and rescale the rest
to a small fraction of their baseline. Linear Green response alone
gives smooth dynamics with cascade $\alpha\approx 1.09$ — no
avalanches. Adding bounded-top-$K$ at $k=12$ drives $\alpha$ into the
SOC band $(2.0, 3.5)$ with $R^{2}>0.85$.

\textbf{Why $k=12$.} The choice $k=12$ is the average degree of
$G_{V_{600}}$ (\S\ref{ssec:graph}), pinned by the substrate
geometry, not by neural data. Smaller $k$ (e.g.\ $k=6$) gives $\alpha$
at the band edge with poorer fit; larger $k$ ($24, 48$) gives $\alpha$
above band or with degraded fit. We do not search $k$ over a fitted
window; $k$ is determined by the graph.

\subsection{The integrated-information proxy
            \texorpdfstring{$\Phi$}{Phi}}

Implementation:
\texttt{kernel/consciousness\_binding.py:phi\_iit\_trajectory}.
Given the state history matrix $S\in\mathbb{R}^{T\times 120}$, write
$A = S\cdot V$ for the H$_4$-eigenvector matrix $V$ (mode amplitudes
$A\in\mathbb{R}^{T\times 120}$). Define $c_{\mathrm{full}}$ as the
lag-$1$ auto-correlation of the full system, and $c_{k}$ as the
lag-$1$ auto-correlation within the K-class irrep block $k$. Then
\[
\Phi \;=\; \max\!\bigl(0,\; |c_{\mathrm{full}}| - \mathrm{mean}_{k}\,|c_{k}|\bigr).
\]
The proxy is designed to be small under H$_{4}$-equivariant dynamics
(when block autocorrelations within irrep classes match the full-system
autocorrelation) and to increase when dynamics produce cross-class
asymmetries. It is not a theorem on information transport; it is a
proxy that captures one observable signature of cross-class
non-equivariance. This is a port of the published
\texttt{integrated\_information\_phi\_irrep} proxy from the cascade
pipeline, adapted to take amplitude trajectories from any source.

We position $\Phi$ as an IIT-style direction-of-effect proxy, not as
a full implementation of IIT. ARIA does not implement cause-effect
repertoires, exclusion postulate, or integration-over-partitions
machinery~\citep{Tononi2008,BalduzziTononi2008}. The propofol $\Phi$
collapse in~\S\ref{sec:results} is consistent with the IIT direction
of effect on the propofol-vs-wake state contrast; it is not a
discharge of the IIT axioms.

\subsection{The continuity composite}

Implementation:
\texttt{kernel/self\_model\_stream.py:StreamContinuityScorer}.
A composite first-person continuity score over a 64-tick rolling
window:
\begin{align*}
b_{\mathrm{cont}} &= \mathrm{mean\ cos\text{-}sim\ of\ consecutive}\ (I, L, P)\,\mathrm{vectors},\\
v_{\mathrm{cont}} &= 1/(1 + 4\cdot \mathrm{var}(\mathrm{valence})),\\
m_{\mathrm{pers}} &= \mathrm{frac.\ of\ consecutive\ same\text{-}modality\ ticks},\\
i_{\mathrm{smooth}} &= 1/(1 + 4\cdot \mathrm{TV}(\mathrm{intensity})),\\
\mathrm{composite} &= 0.35\cdot b_{\mathrm{cont}} + 0.25\cdot v_{\mathrm{cont}} + 0.20\cdot m_{\mathrm{pers}} + 0.20\cdot i_{\mathrm{smooth}}.
\end{align*}
This composite produces the propofol continuity-drop signature
(WAKE composite $0.943$; PROPOFOL composite $0.877$;
drop $+0.066$).

\subsection{The phenomenal-field binding}

Implementation:
\texttt{kernel/consciousness\_binding.py:bind\_phenomenal\_field}.
The substrate state $\psi_{t}$ is mapped to a phenomenal snapshot
with channels (intensity $I$, self-luminosity $L$, presence $P$,
valence, modality\_label). The modality\_label is determined by which
H$_4$ K-class dominates the isotypic compression of $\psi_{t}$ under
the $\sigma$-orbit projector basis. The default profile
\texttt{K\_7\_only} restricts to $\sigma$-twin K-classes for modality
labelling; H$_4$-proper classes contribute amplitude bias.

\subsection{Stimulus models}

Implementation: \texttt{demo\_drug\_sleep\_v4.py}. Four conditions
$\times$ $800$ ticks each at seed $42$:

\textbf{WAKE.} AR(1) cortical noise ($\beta=0.90$), tonic equator-shell
coherence (small always-on bias), and attention episodes (20--50
ticks at amplitude $0.8$, anchored to the largest shell with $15\%$
within-shell rotation per tick). The AR(1) gives temporal correlation
that lets the $\eta=0.20$ self-loop integrate; tonic coherence anchors
modality; attention episodes mimic biological visual fixation
(200--400~ms dwell time analogue); within-shell rotation generates
cascade events without changing modality.

\textbf{SLEEP\_N3.} Slow oscillation ($\sim 1$\,Hz analogue,
amplitude $1.0$) on a coherent shell, plus spindle bursts ($12$ ticks
every $100$ at amplitude $0.4$ fast modulation), plus K-complexes
($4\%$ of ticks at amplitude $0.8$).

\textbf{PROPOFOL.} Low-amplitude tonic noise (amplitude $0.05$);
$\eta = 0.00$ (broken recurrence). Residual cortex preserved as
background drive.

\textbf{RECOVERY.} Identical to WAKE — verifies deterministic
repeatability under the WAKE stimulus protocol after exposure to
PROPOFOL (no hidden persistent modification of the substrate state).

The v4 stimulus models were redesigned after diagnostics on the
v3 stimulus models (which produced 4/6 signatures) to use
biologically-motivated stimulus components — AR(1) cortical noise,
attention episodes, slow-wave drive, spindles, K-complexes — at
amplitudes and durations matching published biological time scales.
They are not fitted to subject-level measurements, but they are
condition-specific design choices iterated to close v3 stimulus-model
artefacts (\texttt{CONSCIOUSNESS\_CHAIN\_V4\_SIGNATURES.md} documents
the v3$\to$v4 redesign). The full stimulus code is
\texttt{demo\_drug\_sleep\_v4.py}.

\subsection{Cascade-mechanism ablation matrix}

The cascade dynamics on the substrate use four mechanisms acting on
the pressure field, each ablatable independently. The $2^{4}$
ablation grid is the basis for the preregistered tests P1--P5 and
the C$\times$P stress test in~\S\ref{sec:stress}.

\textbf{$D$ — D$_{4}$ orbit coupling.} The H$_4$ root system contains
five disjoint 24-cells (D$_4$ orbits). $D$ adds a small
(coupling $0.05$) cross-orbit pressure averaging that prevents
cascades from localising to one orbit.
Implementation: \texttt{kernel/dimensional\_monitor.py:300-305}.

\textbf{$C$ — Context rotation (S$^{7}$ observer frames).} The active
observer frame on the S$^{7}$ rung rotates periodically based on
which uncrossed vertices have accumulated pressure aligning with
each frame's preferences. This creates churn in \emph{which}
vertices are uncrossed at any tick.
Implementation: \texttt{kernel/dimensional\_monitor.py:316-318, 823-827}.

\textbf{$P$ — Partial emission.} High-pressure uncrossed vertices
(above threshold but not yet crossed) emit pressure at $30\%$ scale,
saturating at pressure $3.0$. Without this mechanism, only fully-
crossed vertices emit.
Implementation: \texttt{kernel/dimensional\_monitor.py:842-855}.

\textbf{$E$ — Equator compensation.} The H$_3$ shell-4 equator is a
30-vertex icosidodecahedral ring with split degree distribution.
$E$ scales pressure gain by $(\bar d / d_{v})$ so sparse commissural
vertices overcome their connectivity deficit.
Implementation: \texttt{kernel/dimensional\_monitor.py:320-360}.

The four mechanisms' \emph{targets} are geometry-pinned (D$_4$ orbits,
$S^{7}$ rung, equatorial shell); their gains and coupling constants
($D$ at $0.05$, $P$ at $30\%$ scale saturating at pressure $3.0$,
$C$ rotation period, $E$ degree-ratio multiplier) are fixed design
choices reported here, not learned from data. Their causal effects
within the factorial ablation model are reported in~\S\ref{sec:stress}.
